\subsection{Ablation Protocol}
\label{sec:ablations-setup}

\textbf{Dataset preprocessing.}
Following OXtal~\citep{jin2025oxtal}, we construct the dataset from CSD entries
deposited by May 1, 2025. We require three-dimensional coordinates, an
$R$-factor of at most $9\%$, single-crystal X-ray diffraction at ambient
pressure, a non-polymeric structure, and a known space group. We add missing
hydrogens using the CSD API~\citep{sykes2024scripting} and retain unit cells
with at most 512 atoms, including hydrogens. To prevent benchmark leakage, we
remove all entries belonging to a test CSD family or containing an eligible
CSD-provided component SMILES found in the test sets. Within each remaining CSD
family, we deduplicate equivalent structures using pymatgen
StructureMatcher~\citep{ong2013pymatgen} with $\operatorname{ltol}=0.2$,
$\operatorname{stol}=0.3$, and $\operatorname{angle\_tol}=5^\circ$, retaining the entry
with the lowest $R$-factor. The complete preprocessing pipeline and dataset
statistics are provided in \Cref{app:dataset-preprocessing}.

\textbf{Setup.}
To make controlled ablations feasible under a fixed compute budget, we train each
configuration for 300 epochs on structures with at most 300 atoms per unit cell,
saving checkpoints every 50 epochs. We evaluate each checkpoint on a fixed subset
of 8{,}192 validation crystals, generating 30 candidates per target with the
200-step EDM--Heun sampler, and report $\mathrm{Sol}_C$ as defined in
\Cref{sec:structure-generation}. Architecture, training, conditioning, and scaling
studies follow this protocol; inference studies instead vary the sampler and
function-evaluation budget. Within each study, we vary one design choice at a time.
